\chapter{GroundOne: Method and Implementation}
\label{sec:method}

% STATUS: skeleton. Each section lists what it must contain (TODO notes, not thesis text).

\section{Task and question set}
% TODO: One three-way decision per sentence (supported / contradicted / not mentioned), K = scored sentences; binary score and strict/lenient policies derived from the probabilities (see Ch. 4, Sec. task); readout design for the K decisions (batched suffixes vs per-sentence markers, decided by the prototype).

\section{Prompt format and shared prefix}
% TODO: Document + numbered answer as shared prefix, index-based question suffix; single source of truth for training and inference.

\section{Batched single-pass readout}
% TODO: Prefix prefilled once, KV cache replicated K times, right-padded suffix batch, restricted softmax over letters; exactness (|dp| = 0, Phase 0b). Cost model.

\section{Training objective}
% TODO: LoRA on all linear layers; cross-entropy over the allowed letters at the last prompt token; length-bucketed batches; hyper-parameters (YAML) and the LR sweep.

\section{Calibration}
% TODO: Temperature (+ bias) fitted on dev; bound to the model fingerprint.
% TODO (author's request): define ECE here: bin predictions by confidence into M bins (M = 15, equal width); ECE = sum_m |B_m|/n * |acc(B_m) - conf(B_m)|; binary variant on P(supported) vs observed supported rate, top-label variant for three classes; dependence on binning; never reported alone (a constant base-rate predictor has ECE ~ 0), always with Brier, NLL and reliability diagrams; cite Naeini 2015, Guo 2017.

\section{Software}
% TODO: Package, CLI, native and Jev-compatible API, memory guard; where each advantage comes from (mechanism vs training).

